%% ============================================================
%%  §2  Background and Related Work
%% ============================================================

\section{Background and Related Work}
\label{sec:background}

\subsection{Lean 4 as a Verification Platform}

Lean~4~\cite{lean4} is simultaneously a dependently typed programming language
and an interactive theorem prover. The two roles share a common kernel: a small,
trusted type checker that accepts or rejects proof terms. This kernel is the
trust anchor. A theorem is established if and only if the kernel accepts a
proof term of the correct type.

For our purposes, two properties of Lean~4 are critical. First, \emph{kernel
transparency}: definitions marked \texttt{def} (as opposed to \texttt{opaque}
or \texttt{@[extern]}) are unfolded by the kernel during type-checking. This
means that a closed expression like \texttt{evalFun [] return42Fun 5 []}
can be fully reduced by the kernel to its normal form, making proofs by
\texttt{rfl} possible. Second, \emph{definitional equality}: the kernel
checks proof terms by normalisation, so two terms that reduce to the same
normal form are accepted as equal without requiring explicit rewrite steps.

Lean~4's tactic system is implemented in Lean~4 itself, which makes it
extensible. Custom tactics and macros can be written as Lean~4 programs,
compiled, and used immediately within the same file. This facility is
central to the proof automation described in \S\ref{sec:automation}.

\textbf{Mathlib.} The Mathlib library~\cite{mathlib2020} provides an extensive
collection of mathematical results for Lean~4, including the Fibonacci
sequence (\texttt{Nat.fib} and its properties), linear arithmetic tactics
(\texttt{omega}), and algebraic simplification (\texttt{ring}, \texttt{norm\_num}).
Our project imports \texttt{Mathlib.Tactic} for access to these tools;
the Fibonacci case study (\S\ref{sec:fib}) uses
\texttt{Nat.fib\_add\_two} and \texttt{Nat.fib\_mono} from Mathlib.

\subsection{Rust's Ownership Model and LLBC}

Rust enforces an ``aliasing XOR mutability'' discipline via its type
system~\cite{rustOwnership2014}. At any point, a memory location is either
accessible through multiple immutable references, or through exactly one mutable
reference, but not both. This invariant is checked statically by the borrow
checker and eliminates the data races and use-after-free bugs that afflict C and
C++ programs.

The compiler pipeline proceeds: Rust source $\to$ HIR (high-level IR)
$\to$ MIR (mid-level IR) $\to$ LLVM IR. MIR makes borrow operations, moves,
and drops explicit, but its control-flow graph is unstructured (arbitrary jumps,
implicit panic handlers). The Charon framework~\cite{charon2024} consumes MIR
and produces two higher-level IRs: ULLBC (Unstructured LLBC), which simplifies
pattern matching and constants, and LLBC (Low-Level Borrow Calculus), which
reconstructs structured control flow — loops, matches, sequential branches.
LLBC is the target of our deep embedding.

The key property of LLBC for our purposes is that loops are represented as
a \emph{single} construct: \texttt{loop body} executes \texttt{body}
repeatedly until a \texttt{break} statement is encountered. This maps directly
to the recursive case in our fuel-parameterised interpreter. Conditionals are
represented as \texttt{ite cond trueBranch falseBranch}. Assignments are
\texttt{assign dest rvalue cont}, where the continuation \texttt{cont} is the
next statement.

\subsection{The Forward-Backward Paradigm for \texttt{\&mut T}}

Mutable references in Rust present a semantic challenge for functional models.
When a function returns a \texttt{\&'a mut T}, the caller may subsequently
mutate through that reference, and those mutations must propagate back to
the original storage location. Representing this naively requires a heap model
or pointer arithmetic.

The Aeneas translation~\cite{ho2022aeneas} resolves this with the
\emph{forward-backward} paradigm. A function that operates on mutable borrows
is split into: a \emph{forward function} (modelling execution up to the return
point) and a \emph{backward function} (propagating borrow updates back to
inputs when lifetimes end). Both are pure. We implement the same paradigm
via the \texttt{FBPair} type in \S\ref{sec:monad}.

\subsection{Comparison with Existing Rust Verification Systems}

Table~\ref{tab:comparison} summarises the relationship between lean-pl-verify
and the main competing tools. Several distinctions merit comment.

\begin{table*}[t]
\centering
\caption{Comparison of Rust verification approaches.
         ``TCB'' = trusted computing base; ``Lang.'' = source language(s) supported;
         ``Loop'' = support for unbounded loop verification.}
\label{tab:comparison}
\small
\begin{tabular}{@{}llllll@{}}
\toprule
\textbf{Tool} & \textbf{Approach} & \textbf{TCB} & \textbf{Lang.} & \textbf{Loop} & \textbf{Notes} \\
\midrule
Aeneas~\cite{ho2022aeneas}   & Shallow monadic translation & Charon + code generator & Safe Rust & Yes & Loop summaries \\
Verus~\cite{verus2023}        & SMT-backed contracts & Z3 + verifier & Rust (incl.\ unsafe) & Yes & SMT cliff on complex goals \\
Creusot~\cite{creusot2022}   & Why3 deductive & Why3 + SMT & Safe Rust & Yes & Intrinsic specification \\
Kani~\cite{kani2022}         & Bounded model checking & CBMC/SAT & Rust (incl.\ unsafe) & Bounded & No unbounded loops \\
RustBelt~\cite{rustbelt2018} & Semantic (Iris/Coq) & Coq + Iris & Unsafe Rust & Yes & High manual cost \\
\midrule
\textbf{This work}            & Deep embedding (Lean 4) & Lean kernel only & Rust + TypeScript & Yes & Cross-language unification \\
\bottomrule
\end{tabular}
\end{table*}

\textbf{Aeneas} is the closest technical relative. Both Aeneas and our work
translate safe Rust into Lean~4 and prove functional correctness properties.
The key difference is the embedding strategy. Aeneas generates native Lean
functions via a trusted translation pipeline; correctness of the generated code
depends on correctness of the Charon+Aeneas toolchain. Our deep embedding
makes no such assumption: the interpreter is a Lean~4 \texttt{def}, and its
correctness is definitionally visible to the kernel. A second difference is
scope: Aeneas does not target TypeScript, and cross-language unification is
not a goal. A third difference is proof style: Aeneas encourages users to
write proofs about generated Lean functions, whereas we write proofs directly
about the AST.

\textbf{Verus} offers the strongest industrial story among the compared tools,
with real-world verification of distributed systems components. Its reliance on
Z3 introduces the possibility of false negatives (SMT timeouts) and places the
SMT solver in the TCB. Verus does not support TypeScript.

\textbf{F$^\star$}~\cite{fstar2016} deserves mention: it uses Dijkstra monads
to derive verification condition generators from effectful programs, placing
it philosophically close to our use of \RustM. The difference is that F$^\star$
targets its own language; applying it to Rust would require a different extraction
pipeline.

\textbf{LOOM}~\cite{loom2026} is a recent framework for foundational multi-modal
verification using Dijkstra monads in Lean~4. It shares our interest in
language-agnostic specification but does not target Rust or TypeScript programs
via deep AST embeddings; it operates on shallow monadic terms.

\subsection{Deep \vs Shallow Embeddings}

The choice between deep and shallow embeddings in mechanised semantics is
well-studied~\cite{coqArt}. A shallow embedding maps source constructs directly
to host language constructs; proofs about the source are proofs about the host
language, benefiting from all of the host's proof support. A deep embedding
represents source syntax explicitly; proofs require explicitly reasoning about
the interpreter, but the syntax is a first-class object that can be inspected
and transformed.

For our use case, the deep embedding brings one decisive advantage: the Lean
kernel can evaluate the interpreter on closed terms. A ground computation
like \texttt{evalFun [] fibFun 50 [.int 10]} reduces by kernel reduction
to a concrete result, and \texttt{rfl} proves the correctness lemma. With
a shallow embedding, this would require rewriting along the compiled function
body, which is generally less direct.

\subsection{Note on AI-Assisted Development}

This paper was developed with substantial assistance from Claude Code
(Anthropic), an AI coding agent. The research design, theorem statements,
proof strategies, and scientific conclusions are the author's own. The AI
agent's roles included: generating Lean~4 proof scripts (all verified by
\lakeCmd{lake build}), writing the Charon translation pipeline, scaffolding
proof structures from author-specified statements, generating repetitive proof
instances, and drafting paper text under author direction.

Readers should be aware that a significant portion of the prose and code was
AI-generated and subsequently reviewed and corrected by the author. Variations
in writing style between sections reflect the iterative human-AI editing
process. The complete AI disclosure appears in the Acknowledgements.

All 258 theorems are kernel-checked (0 sorry); no theorem was accepted without
\lakeCmd{lake build} verification. The Lean kernel is the sole trust anchor.
